% Generated draft: requires mathematical and source review.
\documentclass[11pt]{article}
\usepackage[margin=25mm]{geometry}
\usepackage{fontspec}
\setmainfont{Latin Modern Roman}
\usepackage{amsmath,amssymb,mathtools,unicode-math}
\setmathfont{Latin Modern Math}
\usepackage{enumitem,needspace,xurl,hyperref}
\setlength{\parindent}{0pt}
\setlength{\parskip}{5pt}
\setlength{\emergencystretch}{3em}
\newcommand{\N}{\mathbb{N}}
\newcommand{\Z}{\mathbb{Z}}
\newcommand{\Q}{\mathbb{Q}}
\newcommand{\R}{\mathbb{R}}
\newcommand{\C}{\mathbb{C}}
\newcommand{\F}{\mathbb{F}}
\newcommand{\A}{\mathbb{A}}
\newcommand{\PP}{\mathbb{P}}
\newcommand{\E}{\mathbb{E}}
\newcommand{\eps}{\varepsilon}
\newcommand{\dd}{\,\mathrm{d}}
\begin{document}
\section*{UnsolvedMath ID: 30006080}
\textbf{Source problem code:} catalogue entry
\subsection{Definitions and mathematical statement}
The supplied record describes an extension problem, but does not specify enough hypotheses or conclusions to determine a single conjecture.

For notation, let \(A\) denote an abelian variety over a base field that remains unspecified in the source, and put \(g=\dim A\). For a closed subvariety \(X\subseteq A\), put \(d=\dim X\) and \(\operatorname{codim}_A X=g-d\). When the closed immersion is smooth, its normal bundle is
\[
N_{X/A}=\bigl(T_A|_X\bigr)/T_X.
\]
A vector bundle \(E\) is ample if the tautological line bundle \(\mathcal O_{\mathbf P(E)}(1)\) is ample, where \(\mathbf P(E)\) parametrizes one-dimensional quotients. A morphism \(f\colon Y\to A\) is finite if, for every affine open \(\operatorname{Spec}R\subseteq A\), its inverse image is \(\operatorname{Spec}B\) with \(B\) a finite \(R\)-module. Finiteness alone does not require surjectivity onto \(A\).

\textbf{Extension problem as supplied.} Extend arithmetic-finiteness and big-monodromy results for high-codimension subvarieties of abelian varieties with ample normal bundle to subvarieties of smaller codimension, and more generally to varieties \(Y\) admitting a finite morphism \(Y\to A\) to an abelian variety. The record does not state the arithmetic-finiteness or monodromy theorems to be extended. In particular, it specifies neither the arithmetic base and fixed data for a finiteness assertion nor the family, representation, or target group for a monodromy assertion. It also does not specify a smaller-codimension range or which smoothness, irreducibility, and positivity hypotheses should remain in the proposed extension. These are unresolved parts of the formulation; no universal theorem over all base fields or all finite morphisms is asserted here.

\textbf{Quoted contextual assertion.} The original excerpt also contains a statement labelled Theorem 4. In the notation above, it states that if \(X\subset A\) is smooth and irreducible, has ample normal bundle, and satisfies \(d<g/2\), then the following are equivalent:
\begin{enumerate}
\item[(a)] \(X\subset A\) is nondivisible and its Tannaka group satisfies \(G^{*}_{X,\omega}\simeq E_6\).
\item[(b)] \(X\) is isomorphic to the Fano surface of lines on a smooth cubic threefold, and the canonical morphism \(\operatorname{Alb}(X)\to A\) is an isogeny.
\end{enumerate}
Here the Fano surface parametrizes lines contained in the cubic threefold; \(\operatorname{Alb}(X)\) denotes the Albanese variety appearing in the excerpt, and an isogeny of abelian varieties is a surjective homomorphism with finite kernel. The source excerpt does not define nondivisibility, the Tannakian construction, \(\omega\), the superscript \(*\), or the precise form of the group denoted by \(E_6\). It also omits the base-field hypotheses and choices underlying the stated canonical Albanese morphism. Consequently, this equivalence is retained as a contextual source assertion with unresolved qualifications, not as a fully specified theorem verified in this draft. The record does not establish whether an extension of this equivalence itself is part of the intended question.
\subsection{Short English statement}
The supplied question asks how arithmetic-finiteness and big-monodromy results extend to smaller-codimension subvarieties of abelian varieties and to varieties admitting finite maps to abelian varieties. The precise existing results, desired conclusions, base field, and retained hypotheses are missing. An adjacent \(E_6\) classification theorem is preserved as context, with its undefined notation and missing qualifications explicitly recorded.
\subsection{Sources}
\begin{itemize}
\item Supplied local UnsolvedMath registry record: permanent ID \(30006080\), title ``Finiteness for Subvarieties of Abelian Varieties,'' exact source problem code \texttt{catalogue entry}. The record belongs to the Oberwolfach Reports Open Problems collection and assigns proposed year 2024; no individual proposer is supplied.
\item Supplied report citation: \emph{Anabelian Geometry and Representations of Fundamental Groups} (2025), Oberwolfach Reports, pp. 2621--2670, DOI \texttt{10.4171/owr/2024/45}. \url{https://doi.org/10.4171/owr/2024/45}. The original excerpt includes the page marker 2663 and a contextual statement labelled Theorem 4. The supplied literature-source label is ``Oberwolfach Report 45/2024.''
\item Supplied literature source: Lawrence--Sawin, ``Shafarevich for hypersurfaces.'' \url{https://annals.math.princeton.edu/2025/202-3/p01}. This is the supplied attribution and reference label; no fuller bibliographic details are provided.
\item Literature assessment dated exactly \texttt{2026-08-22}: the supplied curation labels the problem ``partially solved'' and says that the report proves substantial low-dimensional and ample-normal-bundle cases, that recent work proves finiteness for smooth hypersurfaces in fixed ample classes in suitable abelian varieties, and that the broader smaller-codimension and finite-to-abelian question remains open. This is supplied dated curation, not an independently verified literature assessment; its time-sensitive status and cited claims require independent checking.
\item The registry supplies statement status \texttt{corrected\_verified} and reports that the official Oberwolfach report was checked, the question isolated from surrounding exposition, and its notation repaired for self-contained publication. Those are supplied provenance claims; no source visit or independent verification is claimed here.
\end{itemize}
\end{document}
